
\maketitle

\begin{abstract}
We prove that a coherent family of observations determines a unique probability
measure on the observable $\sigma$-algebra.  The setting is a directed system of
measurable outcome spaces with surjective refinement maps, together with a
compatible family of finitely additive charges on the cylinder algebra.  The
necessary and sufficient condition for $\sigma$-additive extension is
\emph{collective exhaustion} (CE): no level can assign persistent mass to events
the system as a whole sees as empty.

Two independent routes establish the result.  The \emph{Carathéodory route}
assumes $\sigma$-additive marginals and constructs the global measure directly on
the realization space, with no topological assumptions.  The \emph{Stone route}
assumes only finite additivity and derives $\sigma$-additivity from compactness:
the directed Boolean system dualises to an inverse system of compact Hausdorff
spaces whose limit carries a canonical $\sigma$-additive measure; CE ensures this
measure is supported on the image of the sample space.  Under shared hypotheses,
the two constructions coincide.

CE is \emph{irreducible}: not derivable from any finitary or structural condition
on the query system.  The finite-cofinite counterexample and a model-theoretic
argument via \L{}o\'{s}'s theorem locate the boundary precisely.

The Stone space is the canonical compact completion of the observable structure.
Under (re)construction conditions developed in the companion papers, it is
measure-theoretically isomorphic to the state space itself.
\end{abstract}

\begin{quote}\small
\textbf{2020 Mathematics Subject Classification.}
60A10 (primary); 28A60, 06E15, 28A05 (secondary).

\textbf{Keywords.}
observable extension theorem, collective exhaustion, Carathéodory extension,
Stone duality, Boolean algebra, finitely additive measure, $\sigma$-additivity,
directed system, query system.
\end{quote}

% ------------------------------------------------------------------
\section{Introduction}
\label{sec:intro}
% ------------------------------------------------------------------

The programme of which this paper is a part does not begin with a probability
space but with \emph{distinctions}: the acts of observational separation by
which an observer organises their experience of a system.  State space,
$\sigma$-algebra, and measure are not primitives --- they are what coherent
distinctions, pursued to their conclusion, are forced to produce.  This paper
proves the first step: a coherent family of observational distinctions
determines a unique probability measure.  The further steps --- that this
measure carries canonical dynamics, that the delay structure of observation
(re)constructs the state space, and that finite data can certify
(re)construction --- are carried out in the companion papers.

The answer is a single condition: \emph{collective exhaustion} (CE).  A
compatible family of finitely additive charges extends to a $\sigma$-additive
measure if and only if, whenever a decreasing sequence of events vanishes
globally, some level witnesses the vanishing of mass.  CE rules out purely
finitely additive charges, which assign persistent mass to sequences that
vanish in the limit.

\paragraph{Two routes.}
The \emph{Carathéodory route} (\S\ref{sec:cara}) assumes $\sigma$-additive
marginals and constructs the global measure directly on the realization space.
Sequential common refinements compress any countable cylinder family to a single
level, reducing the $\sigma$-subadditivity check to a single marginal.
By the CE theorem (\S\ref{sec:ce}), $\sigma$-additivity of the marginals is
equivalent to CE.

The \emph{Stone route} (\S\ref{sec:stone}) assumes only finite additivity and
derives $\sigma$-additivity from compactness.  Stone duality converts the
directed Boolean system into an inverse system of compact Hausdorff spaces; a
compatible family of charges assembles into a $\sigma$-additive measure on the
inverse limit.  Descending from the Stone space to $\Omega$ requires two
conditions: Discriminability ensures the Stone embedding is injective, and CE
appears as a support condition --- the measure concentrates on the principal
ultrafilters, i.e., on the image of $\Omega$ inside its Stone compactification.

\paragraph{CE is irreducible.}
CE is not derivable from any finitary or structural condition on the query
system.  The finite-cofinite content on $\mathbb{Q}$ satisfies every algebraic
condition while failing CE; a model-theoretic argument via \L{}o\'{s}'s theorem
shows this is not a proof-search failure.  CE is an irreducible commitment about
the completed infinite, analogous to the axiom of infinity in arithmetic.

\paragraph{The Stone space.}
The Stone space $\mathrm{St}(\mathcal{C})$, constructed as a technical device
for measure extension, is the canonical compact Hausdorff completion of $\Omega$
for the observable topology (\S\ref{sec:bridge}).  CE's two faces --- algebraic
convergence and support on $\Omega$ --- are the same condition seen from inside
and outside the algebra.  When an observation function is cyclic for the Koopman
operator, the Stone space is measure-theoretically isomorphic to the state space;
this connection is developed in the companion papers.

\paragraph{Organisation.}
Section~\ref{sec:setup} introduces query systems, cylinder sets, the observable
$\sigma$-algebra, compatible charges, and the three hypotheses (Surjective Evaluation,
Discriminability, CE).  Section~\ref{sec:ce} establishes the CE theorem and its
irreducibility.  Section~\ref{sec:cara} proves the Carathéodory route.
Section~\ref{sec:stone} proves the Stone route.  Section~\ref{sec:bridge}
draws the two routes together and discusses the Stone space.

% ------------------------------------------------------------------
\section{Setup: query systems and charges}
\label{sec:setup}
% ------------------------------------------------------------------

\subsection{Observational structure}

\begin{definition}[Query system]
\label{def:query-system}
A \emph{query system} is a tuple
\[
  (\iota, \leq,
   \{(\mathrm{Out}(i), \mathcal{F}_i)\}_{i \in \iota},
   \{\pi_{ij}\}_{i \leq j},
   \Omega,
   \{\mathrm{eval}_i\}_{i \in \iota})
\]
consisting of:
\begin{itemize}[noitemsep]
  \item a nonempty partially ordered set $(\iota, \leq)$, the \emph{index set};
  \item for each $i \in \iota$, a measurable space $(\mathrm{Out}(i), \mathcal{F}_i)$,
    the \emph{outcome space at level~$i$};
  \item for each $i \leq j$ in $\iota$, a measurable surjection
    $\pi_{ij} : \mathrm{Out}(j) \to \mathrm{Out}(i)$,
    the \emph{refinement map};
  \item a set $\Omega$, the \emph{realization}, together with maps
    $\mathrm{eval}_i : \Omega \to \mathrm{Out}(i)$ for each $i \in \iota$,
    the \emph{evaluation maps},
\end{itemize}
satisfying the \emph{coherence condition}:
$\pi_{ij} \circ \mathrm{eval}_j = \mathrm{eval}_i$ for all $i \leq j$.
\end{definition}

The index set, outcome spaces, and refinement maps constitute the \emph{observational
structure} of the query system; they encode the distinctions the observer can make
and how finer levels coarsen to coarser ones.  The realization $\Omega$ and
evaluation maps record how that structure is instantiated: $\Omega$ is the set of
states consistent with the system, and coherence says evaluation is compatible with
coarsening.

\subsection{Realization and observables}

The observational structure $(\iota, \leq, \{\mathrm{Out}(i)\}, \{\pi_{ij}\})$
determines a canonical realization: the projective limit
\[
  \Omega
  = \varprojlim_{i \in \iota} \mathrm{Out}(i)
  = \Bigl\{
      x \in \prod_{i \in \iota} \mathrm{Out}(i)
      \;\Big|\;
      \pi_{ij}(x_j) = x_i \text{ whenever } i \leq j
    \Bigr\},
\]
with $\mathrm{eval}_i(x) = x_i$.

\begin{definition}[Common coarsenings and common refinements]
The index set $(\iota, \leq)$ admits \emph{common coarsenings} if every pair
$i, j \in \iota$ has a lower bound: there exists $k \leq i$ and $k \leq j$.
It admits \emph{common refinements} if every pair has an upper bound: there
exists $k \geq i$ and $k \geq j$.  It admits \emph{sequential common
refinements} if every countable family $\{i_n\}_{n \geq 1} \subset \iota$
has an upper bound in $\iota$.
\end{definition}

\noindent Sequential common refinements imply common refinements.  Common
coarsenings ensure that cylinder sets form a $\pi$-system (Lemma~\ref{lem:pi-system}
below); sequential common refinements supply the $\sigma$-subadditivity step in the
Carathéodory route.

\begin{definition}[Cylinder sets]
For $i \in \iota$ and $A \in \mathcal{F}_i$, the \emph{level-$i$ cylinder with
base $A$} is
\[
  \mathrm{Cyl}(i, A)
  \;=\; \{\omega \in \Omega : \mathrm{eval}_i(\omega) \in A\}.
\]
The \emph{cylinder generating family} is
\[
  \mathcal{C}
  \;=\; \{\mathrm{Cyl}(i, A) : i \in \iota,\; A \in \mathcal{F}_i\}.
\]
For each $i \in \iota$, let $B_i$ denote the Boolean algebra of subsets of
$\Omega$ generated by $\{\mathrm{Cyl}(i, A) : A \in \mathcal{F}_i\}$.
Then $\mathcal{C} = \bigcup_{i \in \iota} B_i$.

The \emph{observable $\sigma$-algebra} is $\sigma(\mathcal{C})$, the
$\sigma$-algebra on $\Omega$ generated by all cylinder sets.
\end{definition}

The coherence condition gives, for $i \leq j$,
$\mathrm{Cyl}(i, A) = \mathrm{Cyl}(j, \pi_{ij}^{-1}(A))$,
so the same subset of $\Omega$ can be described as a cylinder at any finer level.

\begin{lemma}[Cylinder $\pi$-system]
\label{lem:pi-system}
Assume $(\iota, \leq)$ admits common coarsenings.  Then $\mathcal{C}$ is a
$\pi$-system: it is closed under finite intersections.
\end{lemma}

\begin{proof}
Given $\mathrm{Cyl}(i, A)$ and $\mathrm{Cyl}(j, B)$, let $k$ be a common
lower bound, with $k \leq i$ and $k \leq j$.  Then
\[
  \mathrm{Cyl}(i, A) \cap \mathrm{Cyl}(j, B)
  = \mathrm{Cyl}(k,\; \pi_{ki}^{-1}(A) \cap \pi_{kj}^{-1}(B)),
\]
which is again a cylinder.
\end{proof}

For $i \leq j$, define the \emph{connecting map}
\[
  \varphi_{ij} : B_i \to B_j,
  \qquad
  \varphi_{ij}(\mathrm{Cyl}(i, A)) = \mathrm{Cyl}(j, \pi_{ij}^{-1}(A)).
\]
This is a well-defined Boolean algebra homomorphism.  The family
$\{B_i, \varphi_{ij}\}$ is a directed system of Boolean algebras.

\subsection{Charges and hypotheses}

\begin{definition}[Compatible family]
A \emph{charge} on a Boolean algebra $B$ is a finitely additive function
$\mu : B \to [0, \infty)$ with $\mu(\emptyset) = 0$.  A family of charges
$\{\mu_i\}_{i \in \iota}$ with $\mu_i$ on $B_i$ is \emph{compatible} if
\[
  \mu_j(\varphi_{ij}(E)) = \mu_i(E)
  \qquad \text{for all } i \leq j \text{ and all } E \in B_i.
\]
\end{definition}

\noindent A compatible family also determines a charge on $\mathcal{C}$ by
$\mu(\mathrm{Cyl}(i, A)) := \mu_i(\mathrm{Cyl}(i, A))$; compatibility
ensures this is well-defined.

The main theorems require three conditions.

\begin{definition}[Surjective Evaluation]
\label{def:eval-surjective}
The query system satisfies \emph{Surjective Evaluation} if
$\mathrm{eval}_i : \Omega \to \mathrm{Out}(i)$ is surjective for every
$i \in \iota$.
\end{definition}

\begin{definition}[Discriminability]
\label{def:discriminability}
The query system \emph{discriminates points} if for every $\omega \neq \omega'$
in $\Omega$ there exists $E \in \mathcal{C}$ with $\omega \in E$ and
$\omega' \notin E$.
\end{definition}

\begin{definition}[Collective Exhaustion]
\label{def:ce}
A charge $\mu$ on $\mathcal{C}$ (the common extension of a compatible family
$\{\mu_i\}$) satisfies \emph{collective exhaustion} (CE) if: whenever
$E_1 \supseteq E_2 \supseteq \cdots$ in $\mathcal{C}$ and
$\bigcap_n E_n = \emptyset$, we have $\mu(E_n) \to 0$.
\end{definition}

\noindent Surjective Evaluation is a structural condition ensuring that every outcome at
every level is realized by some element of $\Omega$; it implies that the
refinement maps $\pi_{ij}$ are surjective (by coherence) and that the connecting
maps $\varphi_{ij}$ are injective.  Discriminability ensures the Stone embedding
$\omega \mapsto \{\text{principal ultrafilter at } \omega\}$ is injective.  CE
is the irreducible valuation-layer condition studied in
\S\ref{sec:ce}; it rules out purely finitely additive charges and is the
necessary and sufficient condition for $\sigma$-additive extensibility.

% ------------------------------------------------------------------
\section{The CE theorem}
\label{sec:ce}
% ------------------------------------------------------------------

The argument of \S\S\ref{sec:ce}--\ref{sec:stone} has a single subject: the
single condition under which a coherent family of observations determines a
probability measure.  It appears in four forms --- as a local continuity criterion
(\S\ref{sec:ce}, single algebra), as a global equivalence across the directed
system (\S\ref{sec:sp1}), as a direct measure-theoretic construction
(\S\ref{sec:cara}), and as a topological derivation via compactness
(\S\ref{sec:stone}).  These are not separate results but four projections of a
single constraint; \S\ref{sec:bridge} identifies the point where all four
coincide.

This section establishes collective exhaustion as the necessary and sufficient
condition for $\sigma$-additive extensibility.  The analysis works at the level
of an abstract directed system of Boolean algebras; it does not require the full
query-system structure of \S\ref{sec:setup}.  To keep this generality clean, we
write $\mathcal{E}_i$ for a generic Boolean algebra at level $i$ and $\ell_i$
for the charge on it.  The connection to query systems is: $\mathcal{E}_i = B_i$
(the Boolean algebra of cylinder sets at level $i$), $O_i = \mathrm{Out}(i)$,
and $\ell_i = \mu_i$.  The full query system structure re-enters in
\S\ref{sec:cara} and \S\ref{sec:stone}.

\subsection{Single-algebra extension}

\begin{definition}[Boolean charge space]
A \emph{Boolean charge space} is a pair $(\mathcal{E}, \ell)$ where $\mathcal{E}$
is a Boolean algebra of subsets of a set $O$ and
$\ell : \mathcal{E} \to [0,1]$ is a normalized finitely-additive valuation.
No $\sigma$-algebra and no topology are assumed on $O$.
\end{definition}

Under Stone's representation theorem~\cite{Stone1936}, $\mathcal{E}$ is
isomorphic to the clopen algebra of a compact totally-disconnected Hausdorff
space $S(\mathcal{E})$, and $\sigma(\mathcal{E})$ is the Baire $\sigma$-algebra
of $S(\mathcal{E})$.  The \emph{gap} $\sigma(\mathcal{E}) \setminus \mathcal{E}$
consists of the Baire sets that are not clopen: events that countable operations
can reach but that no single observable distinction can name.

\begin{theorem}[Extension criterion]
\label{thm:extension-criterion}
A normalized finitely-additive $\ell : \mathcal{E} \to [0,1]$ extends to a
$\sigma$-additive measure on $\sigma(\mathcal{E})$ if and only if $\ell$ is
continuous at $\emptyset$: for every decreasing sequence
$E_1 \supseteq E_2 \supseteq \cdots$ in $\mathcal{E}$ with
$\bigcap_n E_n = \emptyset$ in $\sigma(\mathcal{E})$, we have
$\ell(E_n) \to 0$.  The extension is then unique.
\end{theorem}

\begin{proof}
Standard; see Halmos~\cite{Halmos1950} or Fremlin~\cite{Fremlin2003}
Vol.~1.  Continuity at $\emptyset$ is used in the Carathéodory extension
argument to pass from finite additivity to $\sigma$-additivity on
$\sigma(\mathcal{E})$.
\end{proof}

\begin{remark}[The gap and the regress]
\label{rem:regress}
The condition in Theorem~\ref{thm:extension-criterion} is not verifiable from
within $\mathcal{E}$ alone: confirming $\bigcap_n E_n = \emptyset$ requires
access to $\sigma(\mathcal{E})$.  In a directed system, one might hope a finer
level $j \geq i$ can witness the emptiness that level $i$ cannot see --- but
the obstruction reappears: finite additivity at $j$ does not force continuity
at $\emptyset$ for sequences in $\sigma(\mathcal{E}_j)$.  Something external
to any single level is required, and \S\ref{sec:incompleteness} shows the
directed-system structure alone does not supply it.
\end{remark}

\subsection{Nontrivial refinement}
\label{sec:incompleteness}

\begin{definition}[Nontrivial refinement]
Let $\pi : O_j \to O_i$ be a surjective refinement map, inducing the Boolean
algebra homomorphism $\pi^{-1} : \mathcal{E}_i \to \mathcal{E}_j$.  The
refinement is \emph{nontrivial} if there exists $A \in \mathcal{E}_j$ with
$A \notin \pi^{-1}(\mathcal{E}_i)$.
\end{definition}

\begin{proposition}[Nontrivial refinement introduces new events]
\label{prop:nontrivial}
If $j \geq i$ is a nontrivial refinement, then $\pi^{-1}(\mathcal{E}_i)
\subsetneq \mathcal{E}_j$.  Equivalently, a level $i$ at which
$\mathcal{E}_j = \pi^{-1}(\mathcal{E}_i)$ for every $j \geq i$ is one above
which the refinement order is algebraically trivial: no finer level adds new
discriminative power.
\end{proposition}

\begin{proof}
Nontriviality asserts the existence of $A \in \mathcal{E}_j$ with
$A \notin \pi^{-1}(\mathcal{E}_i)$, directly giving the strict inclusion.
The contrapositive is the stated equivalence.
\end{proof}

\begin{proposition}[Gap nonemptiness]
\label{prop:gap-nonempty}
Let $\mathcal{E}$ be a Boolean algebra of subsets of a countably infinite set
$O$ that is not a $\sigma$-algebra.  Then
$\sigma(\mathcal{E}) \setminus \mathcal{E} \neq \emptyset$.
\end{proposition}

\begin{proof}
If $\mathcal{E}$ is not a $\sigma$-algebra, there exists a countable family
$\{A_n\} \subset \mathcal{E}$ with $\bigcup_n A_n \notin \mathcal{E}$.  By
definition $\bigcup_n A_n \in \sigma(\mathcal{E})$, so it lies in the gap.
\end{proof}

\subsection{CE characterises $\sigma$-additivity}
\label{sec:sp1}

We now state the directed-system version of CE and prove the main equivalence.

\begin{definition}[Collectively exhaustive family]
\label{def:collectively-exhaustive}
A compatible family $\{\ell_i\}$ of normalized finitely-additive charges on a
directed system is \emph{collectively exhaustive} if: for every $i \in \mathcal{I}$
and every decreasing sequence $E_1 \supseteq E_2 \supseteq \cdots$ in
$\mathcal{E}_i$ with $\bigcap_n E_n = \emptyset$, there exists $j \geq i$
such that $\ell_j(\pi_{ij}^{-1}(E_n)) \to 0$.
\end{definition}

\noindent Here $\pi_{ij} : O_j \to O_i$ is the refinement map ($j$ finer than
$i$), so $\pi_{ij}^{-1}(E_n) \subseteq O_j$ is the preimage of $E_n \subseteq O_i$
at the finer level.

\noindent CE is a \emph{valuation-layer} condition: it places a requirement on
how the charges behave, not on the Boolean-algebraic index structure.  It says
that no level can assign persistent mass to events that the system as a whole
sees as empty.

The following proposition establishes that the index-layer structure of a
directed system — sequential upper-directedness and charge compatibility — does
not on its own force $\sigma$-additivity.

\begin{proposition}[Index-layer conditions do not force $\sigma$-additivity]
\label{prop:independence}
There exists a compatible family of normalized finitely-additive charges on a
sequentially upper-directed directed system that fails $\sigma$-additive
extension at every level.
\end{proposition}

\begin{proof}
Let $\iota = \mathbb{N} \cup \{\omega\}$ with $n \leq \omega$ for all
$n \in \mathbb{N}$, directed by this order.  Take all outcome spaces to be
$\mathbb{Q}$, all event algebras to be the finite-cofinite algebra
$\mathcal{E} = \{E \subseteq \mathbb{Q} : E \text{ is finite or }
E^c \text{ is finite}\}$, all refinement maps to be the identity, and
$\ell(E) = 0$ if $E$ is finite, $\ell(E) = 1$ if $E$ is cofinite.
This is a normalized finitely-additive content and compatibility holds trivially.

Fix an enumeration $q : \mathbb{N} \to \mathbb{Q}$ and let
$F_n = \{q(0), \ldots, q(n)\}$.  Each $F_n$ is finite so $\ell(F_n) = 0$,
yet $\bigcup_n F_n = \mathbb{Q}$ and $\ell(\mathbb{Q}) = 1$.
$\sigma$-subadditivity would give $1 \leq \sum_n \ell(F_n) = 0$: contradiction.
\end{proof}

\begin{theorem}[CE Characterisation Theorem]
\label{thm:sp1}
Let $\{\ell_i\}$ be a normalized compatible family of finitely-additive charges
on a directed system indexed by $\mathcal{I}$.  The following are equivalent:
\begin{enumerate}[label=(\roman*)]
  \item The family is collectively exhaustive.
  \item $\ell_i$ extends uniquely to a $\sigma$-additive measure on
        $\sigma(\mathcal{E}_i)$ for every $i \in \mathcal{I}$.
\end{enumerate}
\end{theorem}

\begin{proof}
\emph{(i) $\Rightarrow$ (ii).}
Fix $i \in \mathcal{I}$ and a decreasing sequence $E_n \searrow \emptyset$ in
$\mathcal{E}_i$.  By CE, find $j \geq i$ with $\ell_j(\pi_{ij}^{-1}(E_n)) \to 0$.
By compatibility, $\ell_i(E_n) = \ell_j(\pi_{ij}^{-1}(E_n)) \to 0$.  So $\ell_i$
is continuous at $\emptyset$.  Theorem~\ref{thm:extension-criterion} gives the
unique $\sigma$-additive extension.

\emph{(ii) $\Rightarrow$ (i).}
Suppose $\ell_i$ extends to $\mu_i$ for every $i$.  Let
$E_n \searrow \emptyset$ in $\mathcal{E}_i$.  Then
$\ell_i(E_n) = \mu_i(E_n) \to \mu_i(\emptyset) = 0$ by $\sigma$-additivity
of $\mu_i$ and normalization.  Take $j = i$.
\end{proof}

\begin{remark}[Dissolution of the index-valuation conflation]
\label{rem:dissolution}
The question ``do the structural conditions on a directed system force
$\sigma$-additivity?'' is ill-posed: it conflates conditions on the index
structure with conditions on the charges.  Proposition~\ref{prop:independence}
separates them; Theorem~\ref{thm:sp1} identifies the exact valuation-layer
condition.  Sequential upper-directedness enters only when assembling per-level
extensions into a global measure (\S\ref{sec:cara}), not at the per-level step.
\end{remark}

\subsection{CE is irreducible}
\label{sec:ce-irred}

\begin{proposition}[CE irreducibility]
\label{prop:ce-irred}
Collective exhaustion is not implied by any structural condition on a query
system.
\begin{enumerate}[label=(\roman*)]
  \item \emph{(Particular case.)}  Sequential upper-directedness, compatibility,
        and normalization do not imply CE.
  \item \emph{(Universal case.)}  No condition $\Phi$ expressible in the
        first-order language of query systems implies CE.
\end{enumerate}
\end{proposition}

\begin{proof}
The particular case is Proposition~\ref{prop:independence}: the finite-cofinite
content satisfies all named structural hypotheses and fails CE.

For the universal case, we proceed in three steps.

\textit{Step 1.}
A first-order condition on query systems is preserved under ultraproducts by
\L{}o\'{s}'s theorem~\cite{Los1955}: if $\Phi$ holds in every member of a family
of query systems, it holds in their ultraproduct.

\textit{Step 2.}
Let $\mathcal{U}$ be a nonprincipal ultrafilter on $\mathbb{N}$ extending the
cofinite filter, and let $\delta_n$ be the point mass at $q(n) \in \mathbb{Q}$
(using the enumeration from Proposition~\ref{prop:independence}).  Each $\delta_n$
is a normalized $\sigma$-additive content satisfying every structural condition
including CE.  The ultraproduct $\prod_{\mathcal{U}} \delta_n$ assigns to
$E \subseteq \mathbb{Q}$ the value $\lim_{\mathcal{U}} \mathbf{1}_{q(n) \in E}$,
which is $1$ on cofinite sets and $0$ on finite sets --- exactly the
finite-cofinite content $\ell$ from Proposition~\ref{prop:independence}.

\textit{Step 3.}
By Step~1, any universally valid first-order $\Phi$ holds for $\ell$.  But
$\ell$ fails CE.  Therefore no first-order $\Phi$ can imply CE.

The Yosida--Hewitt decomposition~\cite{YosidaHewitt1952} explains why this
obstruction is permanent.  Every finitely-additive probability on a Boolean
algebra decomposes uniquely as $\ell = \ell_c + \ell_p$, where $\ell_c$ is
$\sigma$-additive and $\ell_p$ is \emph{purely finitely additive}: it
satisfies every algebraic condition while assigning zero to every singleton.
The finite-cofinite content $\ell$ is purely finitely additive ($\ell_c = 0$).
CE is exactly the condition that rules out a nonzero purely finitely additive
part: $\ell_p = 0$.
\end{proof}

\begin{remark}[CE as irreducible primitive]
\label{rem:ce-primitive}
The gap between structural coherence and $\sigma$-additivity is a proved
boundary, not a proof-search failure.  CE is analogous to the axiom of
infinity: not derivable from the remaining structure, not refutable, but the
commitment without which probability cannot proceed.  The ultrafilter-based
content is the canonical witness: it passes every finite consistency test yet
assigns unit mass to events whose infinite intersection is empty.

CE is the \emph{principle of witnessed vanishing}: a distinction may disappear
globally only if that disappearance is witnessed at some finite level of the
refinement order.  No global collapse without local evidence.
\end{remark}


% ------------------------------------------------------------------
\section{Route 1: The Carathéodory extension}
\label{sec:cara}
% ------------------------------------------------------------------

This section establishes the extension theorem by a direct Carathéodory argument.
This is the third face of the argument: the measure constructed directly from
within the algebra.
The input is a query system satisfying Surjective Evaluation, sequential common
refinements, and common coarsenings, together with a compatible family of
$\sigma$-additive marginals.  By Theorem~\ref{thm:sp1}, the $\sigma$-additivity
of the marginals is equivalent to CE applied at each level.

\subsection{Observational determination}

\begin{proposition}[Observational determination]
\label{thm:obs-determination}
Let the query system admit common coarsenings and let $P$, $P'$ be probability
measures on $(\Omega, \sigma(\mathcal{C}))$.  If
$(\mathrm{eval}_i)_\# P = (\mathrm{eval}_i)_\# P'$ for all $i \in \iota$,
then $P = P'$.
\end{proposition}

\begin{proof}
The equality of pushforward laws implies $P$ and $P'$ agree on every cylinder
event.  By Lemma~\ref{lem:pi-system} the cylinder family $\mathcal{C}$ is a
$\pi$-system generating $\sigma(\mathcal{C})$.  The $\pi$--$\lambda$ theorem
gives $P = P'$.
\end{proof}

\subsection{Observable content}

\begin{lemma}[Finite observational content]
\label{thm:finite-content}
Assume the query system admits common coarsenings and common refinements and
satisfies Surjective Evaluation, and let $\{\nu_i\}_{i \in \iota}$ be an
observationally compatible family with each $\nu_i \in \mathcal{P}(\mathrm{Out}(i))$.
Then the assignment
\[
  \mu_0(\mathrm{Cyl}(i, A)) := \nu_i(A),
  \qquad A \in \mathcal{F}_i,
\]
defines a well-defined finitely additive content on $\mathcal{C}$.
\end{lemma}

\begin{proof}
\textit{Well-definedness.}
Suppose $\mathrm{Cyl}(i, A) = \mathrm{Cyl}(j, B)$.  Use common refinements to
find $k$ with $k \geq i$ and $k \geq j$.  The constraint sets in $\mathrm{Out}(k)$
are $\pi_{ik}^{-1}(A)$ and $\pi_{jk}^{-1}(B)$.  The cylinder equality and
Surjective Evaluation (surjectivity of $\mathrm{eval}_k$) imply these coincide in
$\mathrm{Out}(k)$.  Compatibility then gives
$\nu_i(A) = \nu_k(\pi_{ik}^{-1}(A)) = \nu_k(\pi_{jk}^{-1}(B)) = \nu_j(B)$.

\textit{Finite additivity.}
For disjoint cylinders at possibly different levels, compress to a common
refinement $k$ and use additivity of $\nu_k$.
\end{proof}

\subsection{Carathéodory extension theorem}

\begin{theorem}[Observational extension theorem]
\label{thm:obs-extension}
Suppose:
\begin{enumerate}[label=(\roman*)]
  \item $(\iota, \leq)$ admits common coarsenings and sequential common
        refinements;
  \item the query system satisfies Surjective Evaluation;
  \item the family $\{\nu_i\}_{i \in \iota}$ is observationally compatible
        and each $\nu_i$ is a probability measure on
        $(\mathrm{Out}(i), \mathcal{F}_i)$.
\end{enumerate}
Then there exists a unique probability measure $P$ on
$(\Omega, \sigma(\mathcal{C}))$ such that
$(\mathrm{eval}_i)_\# P = \nu_i$ for every $i \in \iota$.
\end{theorem}

\begin{proof}
\textit{Step 1: Finite content.}
Theorem~\ref{thm:finite-content} gives a well-defined finitely additive content
$\mu_0$ on $\mathcal{C}$.

\textit{Step 2: $\sigma$-subadditivity.}
Let $\mathrm{Cyl}(i, B) \subseteq \bigcup_{n \geq 1} \mathrm{Cyl}(i_n, A_n)$.
Apply sequential common refinements to $\{i, i_1, i_2, \ldots\}$ to find a
single index $k$ with $k \geq i$ and $k \geq i_n$ for all $n$.  Compress
each cylinder to level $k$:
\[
  \mathrm{Cyl}(i, B) = \mathrm{Cyl}(k, E),
  \qquad
  \mathrm{Cyl}(i_n, A_n) = \mathrm{Cyl}(k, E_n),
\]
where $E = \pi_{ik}^{-1}(B)$ and $E_n = \pi_{i_n k}^{-1}(A_n)$.
Surjective Evaluation and the cylinder inclusion in $\Omega$ imply $E \subseteq
\bigcup_n E_n$ in $\mathrm{Out}(k)$.  Therefore
\[
  \mu_0(\mathrm{Cyl}(i, B))
  = \nu_k(E)
  \leq \sum_{n=1}^\infty \nu_k(E_n)
  = \sum_{n=1}^\infty \mu_0(\mathrm{Cyl}(i_n, A_n)),
\]
using $\sigma$-subadditivity of the single measure $\nu_k$.

\textit{Step 3: Carathéodory extension.}
The family $\mathcal{C}$ is a semiring generating $\sigma(\mathcal{C})$ and
$\mu_0$ is a $\sigma$-subadditive finitely additive content on it.
Carathéodory's extension theorem~\cite{Billingsley, Bogachev} applies and yields
$P$ on $(\Omega, \sigma(\mathcal{C}))$.

\textit{Step 4: Marginal recovery and normalization.}
For any $A \in \mathcal{F}_i$,
$P(\mathrm{Cyl}(i, A)) = \mu_0(\mathrm{Cyl}(i, A)) = \nu_i(A)$,
so $(\mathrm{eval}_i)_\# P = \nu_i$.  Setting $A = \mathrm{Out}(i)$ gives
$P(\Omega) = 1$.

\textit{Uniqueness.}
Two probability measures with the same evaluation marginals agree on $\mathcal{C}$,
hence on $\sigma(\mathcal{C})$ by Theorem~\ref{thm:obs-determination}.
\end{proof}

\begin{remark}[Relation to Kolmogorov's extension theorem]
\label{rem:kolmogorov}
Kolmogorov's theorem~\cite{Kolmogorov} constructs a probability measure on $S^T$ from consistent
finite-dimensional distributions.  Theorem~\ref{thm:obs-extension} plays the
same role for query systems: the realization space $\Omega = \varprojlim
\mathrm{Out}(i)$ replaces the product $S^T$, compatible marginals replace
consistent finite-dimensional distributions, and sequential common refinements
replace the countable-cofinality condition.  The key conceptual difference is
that no topology is imposed on the outcome spaces: the measure $P$ is built
directly on the measurable structure of $\Omega$ by Carathéodory, without any
appeal to tightness or compactness.
\end{remark}

\begin{remark}[Prokhorov route]
\label{rem:prokhorov}
A third route to the extension theorem starts from \emph{finitely} additive
marginals and derives $\sigma$-additivity from topological compactness.  Under
\emph{surjective projective uniform tightness} --- a condition requiring that
for every $\varepsilon > 0$ there exist compact sets $K_i \subset \mathrm{Out}(i)$
with $\nu_i(\mathrm{Out}(i) \setminus K_i) < \varepsilon$ and
$\pi_{ij}(K_j) = K_i$ for all $i \leq j$ --- the Prokhorov theorem extends
finitely additive marginals on a countable sequentially upper-directed query
system with Hausdorff outcome spaces to a $\sigma$-additive global measure.
This route assumes topological structure on the outcome spaces and is
complementary to the Carathéodory and Stone routes, which require no topology.
\end{remark}

\begin{remark}[Lean formalization]
\label{rem:lean-cara}
Theorem~\ref{thm:obs-extension} is fully formalized in Lean~4
(\texttt{observational\_extension} in \texttt{QuerySystem.lean}, 0 sorrys).
The hypotheses correspond exactly to the four conditions: \texttt{LowerDirected},
\texttt{SequentiallyUpperDirected}, \texttt{EvalSurjective}, and
\texttt{CompatibleMarginals} with \texttt{IsProbabilityMeasure}.
The Carathéodory step uses \texttt{AddContent.measure} from
Mathlib~\cite{Mathlib}.
\end{remark}

% ------------------------------------------------------------------
\section{Route 2: The Stone extension}
\label{sec:stone}
% ------------------------------------------------------------------

This section gives the second, independent proof of the extension theorem.
This is the fourth face: the same measure derived from outside, via compactness.
The input is weaker: charges need only be finitely additive.  $\sigma$-additivity
is derived from the compactness of the Stone space.

The proof proceeds in four steps.  Step~A identifies the Stone space of the
direct limit with the inverse limit of individual Stone spaces.  The inverse
limit measure step uses Choksi's theorem to produce a $\sigma$-additive measure
on the inverse limit.  Step~B1 identifies the pullback of the Borel algebra of
the Stone space with the observable $\sigma$-algebra.  Step~B2 shows that CE
forces the measure to concentrate on the image of $\Omega$.

\subsection{Step A: Stone space of the direct limit}
\label{sec:stepA}

\begin{lemma}[Injectivity of connecting maps]
\label{lem:phi-inj}
If the query system satisfies Surjective Evaluation, then $\varphi_{ij} : B_i \to B_j$
is an injective Boolean algebra homomorphism for every $i \leq j$.
\end{lemma}

\begin{proof}
Surjective Evaluation and the coherence condition give
$\pi_{ij}(\mathrm{eval}_j(\omega)) = \mathrm{eval}_i(\omega)$ for all $\omega$,
so $\mathrm{im}(\pi_{ij}) \supseteq \mathrm{im}(\mathrm{eval}_i) = \mathrm{Out}(i)$:
the refinement maps are surjective.  A surjective map $\pi_{ij}$ has injective
preimage map $\pi_{ij}^{-1}$ on subsets, so
$\varphi_{ij}(\mathrm{Cyl}(i, A)) = \mathrm{Cyl}(j, \pi_{ij}^{-1}(A))$ is
injective on generators, hence on $B_i$.
\end{proof}

\begin{proposition}[Step A]
\label{prop:stepA}
The Stone space of the direct limit satisfies
$\mathrm{St}\!\bigl(\bigcup_i B_i\bigr) \cong \varprojlim_i \mathrm{St}(B_i)$.
\end{proposition}

\begin{proof}
By Lemma~\ref{lem:phi-inj}, the connecting maps $\varphi_{ij}$ are injective
Boolean algebra homomorphisms.  Under Stone duality, injective Boolean algebra
homomorphisms correspond to surjective continuous maps between Stone spaces.
The Stone space functor $\mathrm{St}$ converts colimits in the category of
Boolean algebras to limits in the category of compact Hausdorff spaces
\cite[II.4]{johnstone1982}.  Applied to $\bigcup_i B_i$ with its injective
system, this gives
$\mathrm{St}\!\bigl(\bigcup_i B_i\bigr) \cong \varprojlim_i \mathrm{St}(B_i)$.
\end{proof}

\subsection{The inverse limit measure}
\label{sec:invlim}

\begin{proposition}[Inverse limit measure]
\label{prop:invlim}
Let $\{\mu_i\}$ be a compatible family of normalised charges on $\{B_i\}$.
Then there exists a unique probability measure $\hat{\mu}$ on
$\mathcal{B}(\varprojlim_i \mathrm{St}(B_i))$ whose pushforward along each
projection $\varprojlim_i \mathrm{St}(B_i) \to \mathrm{St}(B_j)$ agrees with
the Borel measure on $\mathrm{St}(B_j)$ corresponding to $\mu_j$.
\end{proposition}

\begin{proof}
\textit{Step 1: Single-algebra correspondence.}
For each $i$, Stone duality identifies $B_i$ with the clopen algebra
$\mathrm{Clop}(\mathrm{St}(B_i))$.  A charge $\mu_i$ on $B_i$ thus defines
a finitely additive set function on the clopens of the compact Hausdorff space
$\mathrm{St}(B_i)$.  Since $\mathrm{St}(B_i)$ is compact and totally
disconnected, the clopens form a base for the topology and a finitely additive
charge on them extends uniquely to a regular Borel measure; see
\cite{cardona2025}, \S6, or \cite{Halmos1950}, \S\S53--54.  Call this measure
$\hat{\mu}_i$ on $\mathrm{St}(B_i)$.

\textit{Step 2: Compatibility.}
The compatibility condition $\mu_j(\varphi_{ij}(E)) = \mu_i(E)$ translates,
under Stone correspondence, to $\hat{\mu}_i$ being the pushforward of
$\hat{\mu}_j$ along the bonding map $\mathrm{St}(B_j) \to \mathrm{St}(B_i)$.
The family $\{\hat{\mu}_i\}$ is compatible on the inverse system
$\{\mathrm{St}(B_i)\}$.

\textit{Step 3: Extension to the inverse limit.}
Each $\mathrm{St}(B_i)$ is compact Hausdorff.  The inverse limit
$\varprojlim_i \mathrm{St}(B_i)$ is compact Hausdorff.  By
\cite[Theorem~1]{choksi1958}, a compatible family of regular Borel probability
measures on a cofiltered inverse system of compact Hausdorff spaces with
surjective bonding maps extends uniquely to a regular Borel probability measure
on the inverse limit.  This gives $\hat{\mu}$ on
$\mathcal{B}(\varprojlim_i \mathrm{St}(B_i))$.
\end{proof}

\subsection{Step B1: the structural identification}
\label{sec:B1}

Let $\mathrm{pure} : \Omega \to \mathrm{St}(\mathcal{C})$ denote the Stone
embedding, which sends $\omega \in \Omega$ to its principal ultrafilter
$\mathrm{pure}(\omega) = \{E \in \mathcal{C} : \omega \in E\}$.
For $E \in \mathcal{C}$, write $[E] = \{u \in \mathrm{St}(\mathcal{C}) : E \in u\}$
for the corresponding basic clopen in $\mathrm{St}(\mathcal{C})$.

\begin{proposition}[Step B1: structural identification]
\label{prop:B1}
Suppose the query system discriminates points.  Then
\[
  \mathrm{pure}^{-1}\!\bigl(\mathcal{B}_{\mathrm{Ba}}(\mathrm{St}(\mathcal{C}))\bigr)
  \;=\; \sigma(\mathcal{C}),
\]
where $\mathcal{B}_{\mathrm{Ba}}$ denotes the Baire $\sigma$-algebra.
\end{proposition}

\begin{proof}
\textbf{($\supseteq$):}
For any $E = \mathrm{Cyl}(i, A) \in \mathcal{C}$,
$\mathrm{pure}^{-1}([E]) = \{\omega : E \in \mathrm{pure}(\omega)\}
= \{\omega : \omega \in E\} = E$.
Since $[E]$ is clopen, hence Baire, every cylinder set is in
$\mathrm{pure}^{-1}(\mathcal{B}_{\mathrm{Ba}})$, and therefore
$\sigma(\mathcal{C}) \subseteq \mathrm{pure}^{-1}(\mathcal{B}_{\mathrm{Ba}})$.

\textbf{($\subseteq$):}
The Stone space $\mathrm{St}(\mathcal{C})$ is compact and totally disconnected,
so its clopens form a topological basis.  The clopens are exactly the sets
$[E]$ for $E \in \mathcal{C}$ \cite[I.3.2]{johnstone1982}, so the Baire
$\sigma$-algebra of $\mathrm{St}(\mathcal{C})$ is
$\sigma(\{[E] : E \in \mathcal{C}\})$.
For any Baire set $V$, $\mathrm{pure}^{-1}(V) \in
\sigma(\{\mathrm{pure}^{-1}([E]) : E \in \mathcal{C}\}) = \sigma(\mathcal{C})$.

\textbf{Injectivity.}
Discriminability implies $\mathrm{pure}$ is injective: if $\omega \neq \omega'$,
there exists $E \in \mathcal{C}$ with $\omega \in E$ and $\omega' \notin E$,
so $E \in \mathrm{pure}(\omega)$ but $E \notin \mathrm{pure}(\omega')$.
Injectivity ensures the pullback does not conflate distinct events.
\end{proof}

\begin{remark}
On a compact Hausdorff space the Baire $\sigma$-algebra is contained in the
Borel $\sigma$-algebra, with equality when the space is second-countable.
We work throughout with the Baire $\sigma$-algebra on $\mathrm{St}(\mathcal{C})$;
this suffices because the measure $\hat{\mu}$ constructed above is a Baire
measure (it is the unique extension of a charge on the clopen algebra), and
Choksi's theorem applies to Baire measures on compact spaces.
\end{remark}

\subsection{Step B2: CE and the support condition}
\label{sec:B2}

\begin{proposition}[Step B2: support condition]
\label{prop:B2}
Let $\mu$ be a charge on $\mathcal{C}$ satisfying CE.  Then the corresponding
Baire measure $\hat{\mu}$ on $\mathrm{St}(\mathcal{C})$ is supported on
$\mathrm{pure}(\Omega)$, the set of principal ultrafilters.
\end{proposition}

\begin{proof}
By the Yosida--Hewitt theorem \cite{YosidaHewitt1952, rao1983}, every bounded
charge on a Boolean algebra decomposes uniquely as $\mu = \mu_c + \mu_p$, where
$\mu_c$ is countably additive and $\mu_p$ is purely finitely additive.  Under
Stone duality, $\mu_c$ corresponds to mass on the principal ultrafilters
$\mathrm{pure}(\Omega) \subset \mathrm{St}(\mathcal{C})$, while $\mu_p$
corresponds to mass on the non-principal ultrafilters~\cite{rao1983}, Ch.~10.

CE is equivalent to $\mu_p = 0$: CE requires $\mu(E_n) \to 0$ whenever
$E_n \searrow \emptyset$, which is exactly the condition that the purely
finitely additive part vanishes (Theorem~\ref{thm:sp1}).  Therefore
$\hat{\mu}_p = 0$ and $\hat{\mu}$ is supported on $\mathrm{pure}(\Omega)$.
\end{proof}

\subsection{The Stone extension theorem}
\label{sec:assembly}

\begin{theorem}[Stone extension theorem]
\label{thm:stone-main}
Let the query system satisfy Surjective Evaluation, Discriminability, and common
coarsenings, and let $\{\mu_i\}$ be a compatible family of normalised charges
on $\{B_i\}$ satisfying CE.  Then there exists a unique probability measure
$P$ on $(\Omega, \sigma(\mathcal{C}))$ such that
\[
  P(\mathrm{Cyl}(i, A))
  \;=\; \mu_i(\mathrm{Cyl}(i, A))
  \qquad \text{for all } i \in \iota \text{ and } A \in \mathcal{F}_i.
\]
\end{theorem}

\begin{proof}
The proof assembles the preceding propositions:
\[
  \{\mu_i\}
  \;\xrightarrow{\S\ref{sec:invlim},\,\text{Step~1}}\;
  \{\hat{\mu}_i\}
  \;\xrightarrow{\S\ref{sec:invlim},\,\text{Step~3}}\;
  \hat{\mu} \text{ on } \varprojlim_i \mathrm{St}(B_i)
  \;\xrightarrow{\text{B1}}\;
  P \text{ on } \sigma(\mathcal{C}).
\]

\textit{Existence.}
By Proposition~\ref{prop:stepA},
$\mathrm{St}(\mathcal{C}) \cong \varprojlim_i \mathrm{St}(B_i)$.
By Proposition~\ref{prop:invlim}, the compatible family $\{\mu_i\}$ extends to
a Baire probability measure $\hat{\mu}$ on
$\varprojlim_i \mathrm{St}(B_i) \cong \mathrm{St}(\mathcal{C})$.

By Proposition~\ref{prop:B2}, $\hat{\mu}$ is supported on
$\mathrm{pure}(\Omega)$.  Since $\mathrm{pure}$ is injective (Discriminability),
the pushforward $P = (\mathrm{pure})_*\hat{\mu}$ --- the image measure under
$\mathrm{pure}^{-1}$ on the support --- is a well-defined probability measure on
$\Omega$.

By Proposition~\ref{prop:B1}, $\mathrm{pure}^{-1}$ maps Baire sets in
$\mathrm{St}(\mathcal{C})$ to sets in $\sigma(\mathcal{C})$, so $P$ is defined
on $\sigma(\mathcal{C})$.  For any cylinder set $E = \mathrm{Cyl}(i, A)$,
\[
  P(E) \;=\; \hat{\mu}([E]) \;=\; \mu_i(\mathrm{Cyl}(i, A)).
\]

\textit{Uniqueness.}
Two probability measures on $(\Omega, \sigma(\mathcal{C}))$ agreeing on all
cylinder sets are equal by Theorem~\ref{thm:obs-determination}, which requires
common coarsenings.  The Stone theorem therefore assumes common coarsenings
in addition to Surjective Evaluation, Discriminability, and CE.
\end{proof}

\begin{remark}[CE is not bypassed by compactness]
\label{rem:ce-not-bypassed}
It may appear that the compactness of the Stone space eliminates the need for CE.
Compactness does ensure that $\hat{\mu}$ on $\mathrm{St}(\mathcal{C})$ is
$\sigma$-additive.  But $\hat{\mu}$ lives on the Stone space, not on $\Omega$.
The descent from $\mathrm{St}(\mathcal{C})$ to $\Omega$ requires that no mass
escapes to the non-principal ultrafilters, and that is precisely CE
(Proposition~\ref{prop:B2}).  In the Stone picture, CE is not an additional
hypothesis but a \emph{support condition}: the measure concentrates on the image
of the sample space inside its Stone compactification.
\end{remark}

\begin{remark}[Lean formalization]
\label{rem:lean-stone}
The Stone route is formalized in Lean~4 in \texttt{StoneDualityExtension.lean}.
The structural results (Tasks~0$'$-A and~0$'$-C: Stone space, cylinder clopens,
evaluation maps, bonding maps) are fully proved with no sorry obligations.
Two intentional sorrys mark Mathlib gaps: the extension of a clopen-algebra
charge to a regular Borel measure on a compact totally-disconnected space
(Halmos \S\S53--54; Fremlin Vol.~1 \S311E), and Choksi's projective-limit
theorem for compatible measures on cofiltered inverse systems of compact Hausdorff
spaces.  The coincidence of routes (Task~0$'$-E:
\texttt{stone\_agrees\_with\_caratheodory}) is fully proved, as an immediate
consequence of \texttt{observational\_determination}.
\end{remark}

% ------------------------------------------------------------------
\section{The bridge}
\label{sec:bridge}
% ------------------------------------------------------------------

\subsection{Two routes, one theorem}

Under shared hypotheses --- Surjective Evaluation, Discriminability, and
$\sigma$-additive compatible marginals --- both routes produce the same measure.

\begin{proposition}[Coincidence of routes]
\label{prop:coincidence}
Suppose the query system satisfies Surjective Evaluation, Discriminability, and
sequential common refinements and common coarsenings, and let $\{\nu_i\}$ be a
compatible family of probability measures.  Let $P^{\mathrm{C}}$ be the measure
produced by Theorem~\ref{thm:obs-extension} and $P^{\mathrm{S}}$ the measure
produced by Theorem~\ref{thm:stone-main}.  Then $P^{\mathrm{C}} = P^{\mathrm{S}}$.
\end{proposition}

\begin{proof}
Both measures agree on every cylinder set: $P^{\mathrm{C}}(\mathrm{Cyl}(i, A))
= \nu_i(A) = P^{\mathrm{S}}(\mathrm{Cyl}(i, A))$.  By observational determination
(Theorem~\ref{thm:obs-determination}), they are equal on $\sigma(\mathcal{C})$.
\end{proof}

\noindent The two routes are genuinely independent proofs of the same theorem.
The Carathéodory route works directly within measure theory; the Stone route
takes a detour through topology, using compactness of the Stone space to derive
$\sigma$-additivity that the Carathéodory route assumes.  That they produce the
same measure is not obvious from the constructions alone; it is a consequence of
uniqueness (Theorem~\ref{thm:obs-determination}), which is doing real work here.

\subsection{CE's two faces}

CE appears in the two routes in structurally different forms, yet it is the same
condition.

In the Carathéodory route (Route~1), CE appears at the level of the charges: the
per-level marginals are $\sigma$-additive, which by Theorem~\ref{thm:sp1} is
equivalent to CE applied at each level.  Here, CE is an \emph{algebraic} condition:
the charges do not assign persistent mass to events that the observer's refinement
hierarchy sees as empty.  It is a condition on the observer's commitments.

In the Stone route (Route~2), CE appears as a \emph{topological} condition: the
Baire measure $\hat{\mu}$ on the Stone space must be supported on
$\mathrm{pure}(\Omega)$, the image of the sample space inside its compact
compactification.  Mass on non-principal ultrafilters corresponds exactly to the
purely finitely additive part $\mu_p$, and CE $\Leftrightarrow \mu_p = 0$
(Proposition~\ref{prop:B2}).  Here, CE is a \emph{geometric} condition: the
observer's measure does not escape to the ``phantom'' points of the Stone
compactification.

The Yosida--Hewitt decomposition is the bridge between the two faces: it
identifies the purely finitely additive part of $\mu$ with the mass that has
escaped to non-principal ultrafilters.  CE says both things simultaneously:
the algebraic convergence is honest, and the geometric measure is confined to
$\Omega$.

\subsection{The Stone space as compact completion}

The Stone space $\mathrm{St}(\mathcal{C})$ constructed in \S\ref{sec:stone} as
a technical device for measure extension is the canonical compact Hausdorff
completion of the sample space for the observable topology.  The Stone embedding
$\mathrm{pure} : \Omega \to \mathrm{St}(\mathcal{C})$ is dense --- every
nonempty basic clopen $[E]$ meets $\mathrm{pure}(\Omega)$ --- and the Stone
space is uniquely determined by $\mathcal{C}$ up to homeomorphism.  It is not
imposed on $\Omega$ from outside but the compactification already contained in
the Boolean algebra of cylinder sets.

Every point in $\mathrm{St}(\mathcal{C})$ is an ultrafilter: a maximal
consistent assignment of yes/no answers to every observable distinction.
Principal ultrafilters correspond to actual sample-space points; non-principal
ultrafilters are the ideal limit points --- consistent distinction patterns the
refinement system is coherent with, but that no element of $\Omega$ realises.
The Stone space is therefore the \emph{space of all consistent distinction
patterns}: the completion of the observer's world to a compact whole.

CE is the condition that closes the loop.  It requires the measure $P$ on
$\Omega$, when extended to $\mathrm{St}(\mathcal{C})$ via $\mathrm{pure}$, to
remain concentrated on the realised part: no mass escapes to the phantom
ultrafilters.  An observer satisfying CE retains no phantom mass: probability remains
confined to the realised world, however far the algebraic completion extends.  The Stone topology, generated by all cylinder clopens,
is the coarsest topology making every observable distinction continuous; in this
sense the Boolean framework is a first approximation, and effective separation
is already graded by the measure --- a refinement developed in the companion
papers.

For a dynamical system $(X, T, \mu)$ with observation function $h$, the
time-delay cylinder algebras form a directed system of exactly the kind treated
here.  The Stone space of the resulting observable algebra is, when $h$ is
cyclic for the Koopman operator $U_T$, measure-theoretically isomorphic to the
state space $X$ itself.  The compact space introduced here as a tool opens the
possibility that the Stone space of the observable algebra \emph{is} the state
space --- a connection made precise under (re)construction conditions developed
in the companion papers.

\bibliographystyle{plain}
\bibliography{references}
